% SEGMENT OLP-0680-S01
% Part: proof-theory
% Chapter: proof-search
% Section: introduction

\documentclass[../../../include/open-logic-section]{subfiles}

\begin{document}

\olfileid{pt}{ps}{int}

\olsection{அறிமுகம்}

\Log{G3c}, \Log{N1i} போன்ற !!{proof} முறைமைகளின்
அடிகோள்களும் அனுமான விதிகளும், தொடரணிகள் அல்லது
!!{formula}s கொண்ட எந்த மரங்கள் சரியான !!{proof}s
எனக் கணிக்கப்படுகின்றன என்பதை வரையறுக்கின்றன.
அத்தகைய தொடரணி அல்லது !!{formula}s மரம் ஒன்று தரப்பட்டால்,
ஒவ்வொரு படியிலும் பயன்படுத்திய விதி, எடுகோள்களுக்கான
அடையாளங்கள், அவற்றை விடுவிக்கும் அனுமான விதிகள் போன்ற
கூடுதல் தகவல்களும் இருக்கலாம். அடிகோள்கள் மற்றும் விதிகளின்
வரையறைகளைக் கொண்டு அந்த மரம் சரியான !!{proof} தானா
என்று சரிபார்க்கலாம். ஆனால், ஒரு குறிப்பிட்ட தொடரணிக்கோ,
கொடுக்கப்பட்ட எடுகோள்களிலிருந்து ஒரு குறிப்பிட்ட
!!{formula} வுக்கோ !!a{proof} ஐ \emph{கண்டுபிடித்தல்}
வேறொரு, பொதுவாகக் கடினமான சிக்கல். அதுவே
\emph{நிறுவல் தேடல்} என்ற சிக்கல்.

% SEGMENT OLP-0680-S02
நிறுவல் தேடலுக்கு மிக எளிய நடைமுறை ஒன்றுள்ளது.
!!{formula}s ஐ முடிவுறு குறியெழுத்துத் தொகுதியிலிருந்து
உருவாகும் குறித்தொடர்களாகக் கருதலாம். தொடரணிகளையும்,
தொடரணிகள் அல்லது !!{formula}s கொண்ட மரங்களையும்,
!!{formula}s இன் குறித்தொடர்களாகக் கருதலாம்; மரத்தின்
கிளைகளைச் சுட்டச் சிறப்பு அடைப்புக்குறிகள் தேவைப்படலாம்.
ஒரு குறித்தொடர் சரியான !!{proof} ஐக் குறிக்கிறதா என்பதை
அடிகோள்களும் விதிகளும் இயந்திர முறையில் சரிபார்க்க உதவுகின்றன.
எனவே \emph{சாத்தியமான எல்லா} சரியான !!{proof}s ஐயும்
இயந்திர முறையில் உருவாக்கலாம்: !!{proof}s ஐக் குறிக்கும்
குறியெழுத்துத் தொகுதியிலிருந்து எல்லாக் குறித்தொடர்களையும்
உருவாக்கி, ஒவ்வொரு வேட்புத் தொடரும் சரியான !!{proof} ஐக்
குறிக்கிறதா என்று சோதிப்போம். இந்த எளிய நிறுவல் தேடல்
நடைமுறை, கொடுக்கப்பட்ட தொடரணியின் !!{proof} ஒன்றைக்
கண்டுபிடிக்கும் வரை எல்லாச் சரியான !!{proof}s ஐயும்
உருவாக்குவதாகும். கொடுக்கப்பட்ட !!{formula} வுக்கான
நிறுவலைத் தேடினால், அதன் விடுவிக்கப்படாத எடுகோள்கள்
கொடுக்கப்பட்டவற்றில் அடங்க வேண்டும்.

% SEGMENT OLP-0680-S03
மேம்பட்ட முறை, கையால் !!a{proof} ஐக் கண்டுபிடிக்கப்
பயன்படுத்தும் இயல்பான முறையைப் போன்றது: இறுதித்
தொடரணியிலிருந்து தொடங்கி, !!a{formula} ஐத் தேர்ந்தெடுத்து,
ஓர் அனுமான விதியைப் ``பின்திசையில்'' பயன்படுத்தி ஒன்று
அல்லது பல முன்கோள்களை உருவாக்குகிறோம். அந்த முன்கோள்களுக்கு
!!{proof}s இருந்தால், அவற்றை இறுதி அனுமானத்துடன் இணைத்துக்
கொடுக்கப்பட்ட தொடரணியின் !!a{proof} ஐப் பெறலாம்.
இயற்கை வருவித்தல் முறைமையிலும் !!a{proof} ஐத் தேட
இதைப் போன்ற, ஆனால் சிக்கலான முறை பயன்படும்.

% SEGMENT OLP-0680-S04
ஆனால் இந்த முறை நிறுவல் இருந்தால் அதைத் தவறாமல்
கண்டுபிடிக்கும் என்று உறுதி தராது. தவறான விதியையோ
!!{formula}s இன் தவறான வரிசையையோ தேர்ந்தெடுத்தால்
மேற்கொண்டு செல்ல முடியாத இடத்தை அடையலாம். மேலும்
இந்த முறை முழுமையாகக் குறிப்பிடப்படவில்லை: ஒவ்வொரு
கட்டத்திலும் ஒன்றுக்கும் மேற்பட்ட !!{formula} களைத்
தேர்ந்தெடுக்கலாம்; பின்திசையில் பயன்படுத்த ஒன்றுக்கும்
மேற்பட்ட விதிகளும் இருக்கலாம். \emph{நிறுவல் தேடல் படிமுறை}
என்பது முழுவதும் குறிப்பிடப்பட்ட நடைமுறை. !!a{proof}
இருந்தால் அதை அது கண்டுபிடிக்கும்; இதுவே இங்கே
``நிறுவலைத் தவறவிடாதது'' என்பதன் பொருள்.

% SEGMENT OLP-0680-S05
சில !!{proof} முறைமைகளுக்கு மற்றவற்றைவிட எளிய
நிறுவல் தேடல் படிமுறைகள் உண்டு. தொடரணி முறைமைகளில்
!!a{formula} இன் !!{main operator} எது என்பதும்,
அது இடப்புறமா வலப்புறமா வருகிறது என்பதும்,
எந்த விதியைப் பின்திசையில் பயன்படுத்த வேண்டும் என்பதைத்
தீர்மானிக்கின்றன. எடுத்துக்காட்டாக,
$\Gamma \Sequent \Delta, !A \land !B$ இன்
வலப்புறத்தில் உள்ள $!A \land !B$ ஐக் கொண்ட
தொடரணிக்கு !!a{proof} ஐத் தேடினால்,
\RightR{\land} விதியைப் பின்திசையில் பயன்படுத்தி
$\Gamma \Sequent \Delta, !A$ மற்றும்
$\Gamma \Sequent \Delta, !B$ என்ற இரு முன்கோள்களை
உருவாக்க வேண்டும். ஆனால் முறைமை \Log{G1c}
எனில் சுருக்க விதி கூடுதல் தேர்வுகளைத் தருவதால்
தேடல் கடினமாகிறது: \RightR{\Contraction} ஐப்
பின்திசையில் பயன்படுத்தி
$\Gamma \Sequent \Delta, !A \land !B, !A \land !B$
என்ற முன்கோளையும் உருவாக்கலாம். பின்னர்
$\Gamma \Sequent \Delta, !A \land !B, !A \land !B, !A \land !B$
என்ற முன்கோள்; இவ்வாறே தொடரலாம். \Log{G2c}
இல் சிக்கல் இன்னும் அதிகம். \RightR{\land} ஐப்
பின்திசையில் பயன்படுத்துவதற்கும்,
$\Gamma = \Gamma_1 \cup \Gamma_2$ ஆகும்படி
$\Gamma_1$, $\Gamma_2$ ஐயும்,
$\Delta = \Delta_1 \cup \Delta_2$ ஆகும்படி
$\Delta_1$, $\Delta_2$ ஐயும் ஊகித்துத் தேர்ந்தெடுக்க
வேண்டும். அப்போதுதான்
$\Gamma_1 \Sequent \Delta_1, !A$ மற்றும்
$\Gamma_2 \Sequent \Delta_2, !B$ என்ற
முன்கோள்களுடன் முயன்று பார்க்கலாம். தவறான ஊகம்
பின்னர் முட்டுச்சந்துக்கு இட்டுச் செல்லலாம்; அப்போது
தொடக்கத்திற்குத் திரும்பி வேறு $\Gamma_1$,
$\Gamma_2$, $\Delta_1$, $\Delta_2$ ஐத்
தேர்ந்தெடுக்க வேண்டும். !!{formula} என்பது
$!A \lor !B$ எனில், \RightR{\lor} க்கான
இரண்டு வாய்ப்புகளில் ஒன்றைத் தேர்ந்தெடுத்து
$\Gamma \Sequent \Delta, !A$ அல்லது
$\Gamma \Sequent \Delta, !B$ என்ற முன்கோளை
உருவாக்க வேண்டும். தவறான தேர்வு மீண்டும்
பின்வாங்கித் தேட வேண்டிய நிலையைத் தரும்.

% SEGMENT OLP-0680-S06
\Log{G3c} முறைமைக்கு இச்சிக்கல்கள் இல்லை. அதில்
அமைப்பியல் விதிகள் இல்லை; !!{main operator} எது,
!!{formula} எந்தப் பக்கத்தில் வருகிறது என்பவற்றால்
பின்திசையில் பயன்படுத்த வேண்டிய அனுமான விதி
தீர்மானிக்கப்படுகிறது. எனவே \Log{G1c} அல்லது
\Log{G2c} ஐவிட \Log{G3c} நிறுவல் தேடலுக்கு
ஏற்றது. வெற்றிகரமான தேடல் \Log{G3c} இல்
!!a{proof} ஐத் தரும். அந்த !!{proof} ஐ
\Log{G1c} அல்லது \Log{G2c} இன்
நிறுவலாக மாற்றலாம்; தேவையெனில்
\Log{N1c} இன் நிறுவலாகவும் மாற்றலாம்.
ஆகவே \Log{G3c} க்கான நிறுவல் தேடல் படிமுறையும்,
\Log{G3c} !!{proof}s ஐ மற்ற முறைமைகளின்
!!{proof}s ஆக மாற்றும் படிமுறைகளும் இருந்தால்,
அந்த முறைமைகளுக்கான நிறுவல் தேடல் சிக்கலுக்கும்
விடை கிடைக்கும்.

% SEGMENT OLP-0680-S07
நிறுவல் தேடல் படிமுறை நிறுவலைத் தவறவிடாது
என்பதை எவ்வாறு அறிவோம்? படிமுறை !!a{proof} ஐக்
கண்டுபிடிக்கத் தவறினால், அத்தகைய !!{proof}
இருக்கவே முடியாது என்பதைக் காட்ட வேண்டும்.
தவறாக வடிவமைக்கப்பட்ட படிமுறை, நிறுவல் இருந்தும்
அதைக் கண்டுபிடிக்கத் தவறலாம். மேலே சொன்ன
எளிய படிமுறைக்கு இந்தப் பிரச்சினை இல்லை;
அது சாத்தியமான எல்லா !!{proof}s ஐயும் தேடுகிறது.
ஆனால் நிறுவல் தேடல் படிமுறை திறமையாகவும்
குறைந்த முயற்சியிலேயே செயல்படவும் வேண்டும்.

% SEGMENT OLP-0680-S08
நிறுவல் தேடல் படிமுறை நிறுவலைத் தவறவிடாது
என்பதைக் காட்டும் மையக் கருத்து இதுதான்:
படிமுறை தோல்வியுற்றால், அதன் தோல்வி ஏற்பட்ட
முறையிலிருந்தே தொடரணிக்கு !!a{proof} இருக்க
முடியாது என்பதைக் காட்டலாம். அதற்குத்
தொடரணி ஏற்புடையதல்ல என்பதைக் காட்டுவோம்.
செவ்வியல் முதல்நிலைத் தருக்கத்தில்,
$\Gamma \Sequent \Delta$ போன்ற ஏற்புடைய
தொடரணிகளின் பண்பு: ஒவ்வொரு !!{structure} இலும்
$\Gamma$ விலுள்ள குறைந்தது ஒரு !!{formula}
பொய்யாகவோ, $\Delta$ விலுள்ள குறைந்தது ஒரு
!!{formula} மெய்யாகவோ இருக்கும். \Log{G3c}
\emph{சரித்தன்மையுடையது}; அதாவது அது ஏற்புடைய
தொடரணிகளை மட்டுமே !!{prove}s செய்கிறது.
!!{proof}s மீது தொகுத்தறிந்து இதை நிறுவுகிறோம்:
அடிகோள்கள் வெளிப்படையாக ஏற்புடையவை; ஒரு விதியின்
முன்கோள்கள் ஏற்புடையவையானால் அதன் முடிவும் ஏற்புடையது.
நிறுவல் தேடல் படிமுறை நிறுவலைத் தவறவிடாது என்பதைக்
காட்ட, $\Gamma \Sequent \Delta$ க்கான
!!a{proof} ஐத் தேடல் தோல்வியுற்றால்,
$\Gamma$ விலுள்ள எல்லா !!{formula}s உம் மெய்யாகவும்,
$\Delta$ விலுள்ள எல்லா !!{formula}s உம் பொய்யாகவும்
இருக்கும் !!a{structure} ஐ வரையறுக்க முடியும்
என்று காட்டுகிறோம். இதனால் \Log{G3c}
\emph{நிறைவுடையது} என்பதும் நிறுவப்படுகிறது:
$\Gamma \Sequent \Delta$ ஏற்புடையதெனில்,
அதற்கு \Log{G3c}-!!{proof} உண்டு. தோல்வியுற்ற
ஒவ்வொரு நிறுவல் தேடலும்
$\Gamma \Sequent \Delta$ ஏற்புடையதல்ல என்பதைக்
காட்டுவதால், ஏற்புடைய தொடரணிக்கான ஒவ்வொரு
நிறுவல் தேடலும் வெற்றி பெற வேண்டும்.

\end{document}
